%% APPENDIX E -- radio-frequency interference screening.
%% Owner: appendix-triage.  Carries \label{app:rfi}.
\section{Radio-frequency interference screening}
\label{app:rfi}

Interference is the first alternative explanation for an unattributed
crossing. An interferometer is a poor detector of it -- a signal not in the far
field at the phase centre is suppressed in the correlation by its rapidly
varying geometric phase -- though allocations exist across the searched bands
\citep{ITU2024}, and the documented hazard, spaceborne cloud radar at
\RadarLo--\RadarHi\,GHz \citep{CloudSat2002,EarthCARE2023,ALMAMemo504},
overlaps \RadarNWin{} of the \RadarNTot{} census windows, with
\RadarNCross{} crossings and \RadarNStage{} rank-passing events.
\emph{The tests of Table~\ref{tab:rfitests} carry none of the signatures
interference leaves}, but their populations differ, each needing a different
quantity recorded, and since the correlator flagging was not retained
\emph{no per-channel interference flag exists, and none can be recovered}.
The one concentration a reader will notice, \TlBandAll{} crossings between
\TlBandLo{} and \TlBandHi\,GHz, is the carbon monoxide tuning, \TlBandObs{} of
the \TlBandNSub{} unattributed crossings the census and the reserved hold-out
together place in a searched window (row 2.4 of
Table~\ref{tab:counts}) against \TlBandNull$\,\pm\,$\TlBandNullSd{} expected
from where the windows were tuned. The image-sideband test is a bound, not a
clearance, and \ImNDsb{} crossings lie where the receivers are
double-sideband. The full treatment is in
\texttt{documentation/\allowbreak appendix\_E\_rfi.md} of the data archive
\citep{SIdeposit}.

\begin{table}
\centering
\footnotesize
%% ★ v4.26 (r25-captions): 4 printed lines -> the author's 3.  The reason the
%% five populations differ MOVED to the paragraph above, where the record
%% that is or is not retained is the subject.
\caption{The interference tests, each of which could have failed. $p$ is
against the same peaks redrawn on their own channel grids.}
\label{tab:rfitests}
\begin{tabular}{@{}p{0.20\columnwidth}p{0.22\columnwidth}p{0.33\columnwidth}r@{}}
\hline
Test & Population & Result & $p$ \\
\hline
Clustering in sky frequency &
\AppMNCrossF{} census crossings with a released peak &
\AppMTubeObs{} within $\pm\DispoHalfKms$\,km\,s$^{-1}$ of a searched
transition against \AppMTubeNull{} expected: they cluster on the survey's own
lines &
\AppMTubeP \\
Crowding on one another &
\QcNQP{} unattributed, quality-passing, of the census &
no crowding at any of six separations from \QcTolLo{} to \QcTolHi\,MHz, where
the \QcNAttr{} attributed crossings crowd at every one &
$\geq\QcPQPMin$ \\
Atmospheric transitions &
\TlNCross{} crossings of the census and the hold-out, against \TlNLine{}
telluric lines &
\ScNInRad{} inside the \ScRadMHz\,MHz screening radius in the topocentric
frame, flagged as atmospheric and removed from both sides of the chance
comparison while remaining among the unattributed; \TlTubeObs{} inside
$\pm\TlTubeHalf$\,km\,s$^{-1}$ against \TlTubeNull{} expected &
-- \\
Fixed intermediate frequency &
\IfPopNWin{} census windows with a released peak frequency &
peak positions within their own windows uniform, median \RfiFracMed &
-- \\
Cross-target occupancy &
\RfiOccNWin{} extracted windows toward \RfiOccNStars{} stars, every block set
of Table~\ref{tab:blockledger} together &
\RfiOccObs{} cross-star coincidences within \RfiOccTol\,MHz against
\RfiOccNull$\,\pm\,$\RfiOccNullSd{} expected; \emph{below} the null over the
\RfiOccSigN{} above-trigger windows, \RfiOccSigObs{} against
\RfiOccSigNull$\,\pm\,$\RfiOccSigNullSd &
\RfiOccP \\
Image sideband &
the \ImNScreen{} unattributed crossings lying in a catalogued block &
\ImNNoParent{} have no catalogued transition in the image interval and
\ImNParentAbsent{} have one the same block observed with nothing bright
enough to leak; for \ImNLeft{} the parent frequency was not observed &
-- \\
\hline
\end{tabular}
\end{table}



\emph{One execution block fails these tests, as a block and not as a
carrier.} Its \DqBlockNCross{} crossings toward
\LgVerRejectStar{} are the \DqBlockNCross{} highest in the survey, none below
$T_\star=\EvFlagTStarLo$ (\S\ref{sec:bdblock}), and its \BdNSibling{} sibling
executions produce no crossing. \emph{None of this substitutes for the ON/OFF
cadence a dedicated search uses}: these are statements about the population,
not a clearance of any one crossing.

%% "The null result rests on internal validation alone" lives once, in the
%% discussion's lessons subsection, and not here.  fig:appm is not included
%% -- it compared our own two statistics and the paper states only the one it
%% uses -- so figures/appm_power.pdf and the \AppMPowerOld* and
%% \RfiOcc{Prior,Dir,Stale}* families will be retired.
%%
%% ★ r16-mask: the cross-target screen of the fourth test is now scoped to
%% the ADOPTED attribution (round 660, maskrfi_v660.py), which makes it two
%% events and not one.  round 49's \RfiEventLead, \RfiFreqThem,
%% \RfiSameChanClause, \RfiNotTestableClause, \RfiOtherWin,
%% \RfiOtherStarsRange and \RfiOtherBlocksRange are therefore unreferenced
%% and will be retired; rfi_v399.py still reproduces them and its grammar
%% selftest still runs, so nothing breaks if the scope is ever restored.
